\documentclass[10pt,a4paper]{article}

\usepackage[T1]{fontenc}
\usepackage[utf8]{inputenc}
\usepackage[ngerman]{babel}
\usepackage[a4paper,margin=14mm]{geometry}
\usepackage{lmodern}
\usepackage[table]{xcolor}
\usepackage{tabularx}
\usepackage{array}
\usepackage{microtype}
\usepackage{graphicx}
\usepackage[most]{tcolorbox}
\usepackage{tikz}
\usetikzlibrary{positioning,calc}

\pagestyle{empty}
\setlength{\parindent}{0pt}
\setlength{\parskip}{0pt}
\renewcommand{\familydefault}{\sfdefault}
\renewcommand{\arraystretch}{1.16}

\definecolor{Navy}{HTML}{17324D}
\definecolor{Blue}{HTML}{3B6EA5}
\definecolor{Sky}{HTML}{EAF3FA}
\definecolor{Mint}{HTML}{E8F5EF}
\definecolor{Gold}{HTML}{F3B33D}
\definecolor{Ink}{HTML}{1E2935}
\definecolor{Muted}{HTML}{617080}
\definecolor{Line}{HTML}{D8E1E8}
\definecolor{CodeBg}{HTML}{F5F7F9}
\definecolor{KeyBg}{HTML}{FFF3D6}
\definecolor{KeyLine}{HTML}{B87810}

\color{Ink}
\newcommand{\sectiontitle}[1]{%
  \vspace{2.2mm}{\color{Blue}\bfseries\large #1}\par
  \vspace{0.7mm}{\color{Blue}\rule{\linewidth}{0.7pt}}\par\vspace{1.1mm}}
\newcommand{\classcard}[2]{%
  \fcolorbox{Blue}{white}{\parbox{55mm}{\centering\vspace{1mm}{\large\bfseries\ttfamily #1}\par
  \vspace{1mm}{\color{Blue}\rule{\linewidth}{0.6pt}}\par\vspace{1mm}\raggedright\ttfamily\small #2\vspace{1mm}}}}

\begin{document}

\colorbox{Navy}{\parbox{\dimexpr\linewidth-2\fboxsep\relax}{%
  \vspace{3mm}\color{white}\begin{tabularx}{\linewidth}{@{}Xr@{}}
  {\small\bfseries INFORMATIK · KLASSE 10} & \colorbox{Gold}{\color{Navy}\small\bfseries\strut LÖSUNGEN}\end{tabularx}
  \vspace{2.4mm}{\fontsize{18}{21}\selectfont\bfseries Aufgabe 4 -- Klassenkarte und Tabellenschema}\par
  \vspace{0.8mm}{\large Sicherungsblatt}\vspace{2.8mm}}}

\vspace{2.5mm}
\colorbox{Sky}{\parbox{\dimexpr\linewidth-2\fboxsep\relax}{\vspace{1.6mm}
\textcolor{Blue}{\bfseries Ziel:} Du wandelst Tabellenschemata und Klassenkarten um und stellst eine 1:n-Beziehung fachlich korrekt dar.\vspace{1.6mm}}}

\sectiontitle{1 · Zwei Umwandlungsrichtungen}
{\small\color{Muted}\textbf{Klassenkarte → Tabellenschema:} Attribute werden zu Spalten; Datentypen werden ergänzt.}\par\vspace{1.5mm}
\begin{minipage}[t]{0.43\linewidth}
\textcolor{Blue}{\bfseries Tabellenschema}\par\vspace{1mm}
\colorbox{CodeBg}{\parbox{\dimexpr\linewidth-2\fboxsep\relax}{\ttfamily\small
users (\par
\quad id: int\par
\quad username: varchar(255)\par
\quad created\_at: date\par
\quad updated\_at: date\par
)\par\medskip
photos (\par
\quad id: int\par
\colorbox{KeyBg}{\strut\quad user\_id: int}\par
\quad description: varchar(255)\par
\quad url: varchar(255)\par
\quad created\_at: date\par
\quad updated\_at: date\par
)}}
\par\vspace{1mm}
\fcolorbox{KeyLine}{KeyBg}{\strut\scriptsize\textbf{Fremdschlüssel:} \texttt{photos.user\_id}}
\end{minipage}
\hfill
\begin{minipage}[t]{0.54\linewidth}
\centering
\textcolor{Blue}{\bfseries Klassendiagramm}\par\vspace{5mm}
\begin{tikzpicture}[every node/.style={inner sep=0pt}]
\node (users) {\resizebox{38mm}{!}{\classcard{users}{id\par username\par created\_at\par updated\_at}}};
\node[right=25mm of users] (photos) {\resizebox{38mm}{!}{\classcard{photos}{description\par url\par created\_at\par updated\_at}}};
\draw[Blue,line width=.8pt] (users.east) -- node[midway,fill=white,inner sep=1pt,font=\scriptsize] {hochgeladen von} (photos.west);
\node[above=1mm of $(users.east)!0.10!(photos.west)$,font=\bfseries] {1};
\node[above=1mm of $(users.east)!0.90!(photos.west)$,font=\bfseries] {n};
\end{tikzpicture}
\par\vspace{4mm}
{\scriptsize\color{Muted}\textbf{Tabellenschema → Klassendiagramm:} Datentypen entfallen; der Fremdschlüssel wird durch die Beziehung mit 1 und n dargestellt.}
\end{minipage}

\sectiontitle{2 · Die Beziehung lesen}
\begin{tcolorbox}[colback=white,colframe=Line,arc=2mm,boxrule=.7pt,left=2.5mm,right=2.5mm,top=1.5mm,bottom=1.5mm]
\textbf{1:n-Beziehung:} Ein \texttt{user} kann mehrere \texttt{photos} hochladen. Jedes \texttt{photo} gehört zu genau einem \texttt{user}.\par\smallskip
Die Kardinalität \textbf{1} steht bei \texttt{users}; die Kardinalität \textbf{n} steht bei \texttt{photos}.
\end{tcolorbox}

\sectiontitle{3. Merke: 1:n-Beziehung im Tabellenschema und im Klassendiagramm}
\begin{tcolorbox}[colback=Mint,colframe=Blue,arc=2mm,boxrule=.8pt,left=3mm,right=3mm,top=2mm,bottom=2mm]
Der Fremdschlüssel (hier: \texttt{photos.user\_id}) wird nicht als Attribut in die Klassenkarte (hier: \texttt{photos}) übernommen. Im Klassendiagramm zeigt stattdessen die Beziehung (hier: zwischen \texttt{users} und \texttt{photos}) mit \textbf{1} und \textbf{n} diese Information.\par\smallskip
\textbf{Fremdschlüssel:} \texttt{photos.user\_id}\par
Er verweist im Tabellenschema auf \texttt{users.id}.
\end{tcolorbox}

\vfill{\color{Line}\rule{\linewidth}{.5pt}}\par\vspace{1.5mm}
\begin{tabularx}{\linewidth}{@{}Xr@{}}{\small\color{Muted}Klasse 10 · Datenbanken} & {\small\color{Muted}Sicherungsblatt · Aufgabe 4}\end{tabularx}
\end{document}
